\documentclass[12pt,a4paper]{article}

\usepackage[utf8]{inputenc}
\usepackage[T1]{fontenc}
\usepackage[french]{babel}
\usepackage{lmodern}
\usepackage{geometry}
\geometry{top=2.5cm, bottom=2.5cm, left=3cm, right=3cm}
\usepackage{parskip}
\usepackage{xcolor}
\usepackage{booktabs}
\usepackage{tabularx}
\usepackage{float}
\usepackage{fancyhdr}
\usepackage{hyperref}
\usepackage{listings}

% Couleurs
\definecolor{myblue}{RGB}{37,99,235}
\definecolor{mygray}{RGB}{100,116,139}
\definecolor{mylight}{RGB}{241,245,249}
\definecolor{codebg}{RGB}{248,250,252}
\definecolor{codekey}{RGB}{99,102,241}
\definecolor{codestr}{RGB}{16,185,129}
\definecolor{codecmt}{RGB}{148,163,184}

% Style listings -- langue neutre par defaut
\lstset{
  backgroundcolor=\color{codebg},
  basicstyle=\ttfamily\small,
  keywordstyle=\color{codekey}\bfseries,
  stringstyle=\color{codestr},
  commentstyle=\color{codecmt}\itshape,
  breaklines=true,
  frame=single,
  rulecolor=\color{myblue!40},
  numbers=left,
  numberstyle=\tiny\color{mygray},
  numbersep=8pt,
  showstringspaces=false,
  tabsize=2,
  xleftmargin=12pt,
  aboveskip=6pt,
  belowskip=6pt,
}

% En-tete
\setlength{\headheight}{14pt}
\pagestyle{fancy}
\fancyhf{}
\fancyhead[L]{\small\textcolor{mygray}{PET Platform -- Rapport de developpement}}
\fancyhead[R]{\small\textbf{\textcolor{myblue}{Issues \#6 \& \#7}}}
\fancyfoot[C]{\small\thepage}
\fancyfoot[L]{\small\textit{Abdelmoumen Elimrany}}
\fancyfoot[R]{\small 22 avril 2026}
\renewcommand{\headrulewidth}{0.4pt}

% Hyperliens
\hypersetup{
  colorlinks=true, linkcolor=myblue, urlcolor=myblue,
  pdftitle={PET Platform - Rapport}, pdfauthor={Abdelmoumen Elimrany}
}

% Commande ok
\newcommand{\ok}{\textcolor{green!60!black}{\textbf{OK}}}
\newcommand{\statut}[1]{\texttt{\textcolor{myblue}{#1}}}

% ============================================================
\begin{document}

% Page de titre
\begin{titlepage}
\thispagestyle{empty}
\vspace*{3cm}
\begin{center}
  {\color{myblue}\rule{\linewidth}{2pt}}\\[0.5cm]
  {\Huge\bfseries PET Platform}\\[0.3cm]
  {\large\color{mygray} Privacy-Enhancing Technologies}\\[0.2cm]
  {\color{myblue}\rule{\linewidth}{2pt}}\\[1.5cm]

  {\LARGE\bfseries Rapport de developpement}\\[0.4cm]
  {\large Implementation des issues \#6 et \#7}\\[2cm]

  \begin{tabular}{ll}
    \textbf{Etudiant}  & Abdelmoumen Elimrany \\[4pt]
    \textbf{Encadrant} & Achahbi \\[4pt]
    \textbf{Date}      & 22 avril 2026 \\[4pt]
    \textbf{Depot}     & \url{github.com/ElimranyAbdelmoumen/PET\_Projet} \\[4pt]
    \textbf{Branche}   & \texttt{main} (commit \texttt{bb816b0}) \\
  \end{tabular}
\end{center}
\vfill
\begin{center}
  \colorbox{myblue}{\textcolor{white}{\textbf{~Issue \#6~}}}~Support R, Microdata, Collecte des outputs
  \quad
  \colorbox{black}{\textcolor{white}{\textbf{~Issue \#7~}}}~Navigateur securise Brave
\end{center}
\end{titlepage}

% Resume
\newpage
\thispagestyle{fancy}

\noindent\colorbox{mylight}{%
  \begin{minipage}{\linewidth}
    \vspace{8pt}
    \textbf{\large Resume executif}\\[6pt]
    Ce rapport presente l'implementation de deux issues sur la \textbf{PET Platform},
    une application web Flask de soumission securisee de scripts scientifiques Python/R.\\[4pt]
    \textbf{Issue \#6} ajoute : execution R (\texttt{Rscript --vanilla}), selecteur de fichier
    microdata, et workflow de collecte/validation des fichiers de sortie.\\[4pt]
    \textbf{Issue \#7} securise l'acces via l'extension Brave Security Shield (MV3) :
    protections anti-exfiltration et enforcement HTTP 403 cote Flask.\\[4pt]
    Le \textbf{benchmark} (8 scenarios x 5 runs) confirme une latence moyenne de
    \textbf{1,5 secondes} sans aucune erreur sur 40 executions.
    \vspace{8pt}
  \end{minipage}%
}

\vspace{0.5cm}
\tableofcontents
\newpage

% ============================================================
\section{Presentation du projet}

\subsection{Contexte}

La \textbf{PET Platform} permet a des chercheurs de soumettre des scripts de traitement
de donnees pour execution dans un environnement \textbf{entierement isole}. Un
administrateur valide les resultats avant qu'ils soient accessibles a l'utilisateur,
garantissant que les donnees microdata sensibles ne quittent jamais le serveur.

\subsection{Architecture technique}

\begin{table}[H]
\centering
\caption{Composants de la PET Platform}
\begin{tabularx}{\textwidth}{@{}>{\bfseries}lllX@{}}
\toprule
Composant & Technologie & Version & Role \\
\midrule
Backend       & Flask + Jinja2    & Python 3.11 & Interface web, auth, routes \\
Base donnees  & PostgreSQL        & 16          & Soumissions, outputs, users \\
Worker        & Python            & 3.11        & Execution asynchrone \\
Sandbox       & Docker isole      & --          & Conteneur ephemere par soumission \\
Navigateur    & Brave + Extension & MV3         & Acces securise, protections client \\
\bottomrule
\end{tabularx}
\end{table}

\subsection{Cycle de vie d'une soumission}

\begin{table}[H]
\centering
\caption{Etats d'une soumission}
\begin{tabularx}{\textwidth}{@{}clX@{}}
\toprule
Etape & Etat & Description \\
\midrule
1 & \statut{PENDING}  & L'utilisateur soumet le script \\
2 & \statut{APPROVED} & L'administrateur approuve l'execution \\
3 & \statut{RUNNING}  & Le worker lance le conteneur Docker \\
4 & \statut{FINISHED} & Execution terminee, outputs collectes \\
5 & \statut{REJECTED} & L'administrateur peut rejeter a l'etape 2 \\
\bottomrule
\end{tabularx}
\end{table}

% ============================================================
\section{Issue \#6 -- Support R, microdata et collecte des outputs}

\noindent\colorbox{myblue!10}{\parbox{\linewidth}{\vspace{4pt}
\textbf{Objectif :} (1) Executer des scripts R en plus de Python,
(2) associer un fichier microdata a chaque soumission,
(3) collecter et valider les fichiers de sortie produits.
\vspace{4pt}}}

\vspace{6pt}

\subsection{Support R dans le runner Docker}

\subsubsection*{Dockerfile -- ajout de R}

\begin{lstlisting}
# Installation de R
RUN apt-get update && apt-get install -y r-base r-base-dev \
    && rm -rf /var/lib/apt/lists/*

# Packages Python + rpy2
RUN pip install --no-cache-dir pandas numpy matplotlib \
    scikit-learn scipy seaborn rpy2

# Repertoires de sortie (droits ouverts)
RUN mkdir -p /work/outputs /work/data && chmod 777 /work/outputs /work/data

ENTRYPOINT ["python", "/usr/local/bin/run_script.py"]
\end{lstlisting}

\textbf{Correction permissions :} L'utilisateur \texttt{runner} (UID 10001) a ete supprime
du Dockerfile -- il empechait l'ecriture dans \texttt{/work/outputs/} car le repertoire
hote etait detenu par root.

\subsubsection*{Detection Python / R dans le worker}

\begin{lstlisting}
# worker.py -- selection de l'interpreteur selon l'extension
ext = os.path.splitext(file_path)[1].lower() or ".py"
container_script = f"/work/script{ext}"
# .py -> exec() Python avec capture des variables
# .R  -> Rscript --vanilla /work/script.R
\end{lstlisting}

\begin{lstlisting}
# run_script.py -- routing Python / R
if script_path.endswith(".R"):
    result = subprocess.run(
        ["Rscript", "--vanilla", script_path],
        capture_output=True, text=True, timeout=55
    )
else:
    # exec() avec capture automatique (result, output, df)
    ...
\end{lstlisting}

\subsection{Selecteur de fichier microdata}

Les fichiers microdata (CSV, Excel) sont stockes en base avec un UUID.
L'utilisateur selectionne le fichier dans un menu deroulant lors de la soumission.

\begin{lstlisting}
-- Schema base de donnees (db/init.sql)
CREATE TABLE microdata_files (
    id         SERIAL PRIMARY KEY,
    guid       UUID DEFAULT gen_random_uuid() UNIQUE,
    name       VARCHAR(255) NOT NULL,
    file_path  TEXT NOT NULL,
    created_at TIMESTAMPTZ DEFAULT NOW()
);

ALTER TABLE submissions
    ADD COLUMN microdata_guid UUID REFERENCES microdata_files(guid),
    ADD COLUMN outputs_status VARCHAR(20) DEFAULT 'NONE';
\end{lstlisting}

Le fichier est monte \textbf{en lecture seule} dans le conteneur sous
\texttt{/work/data/<nom\_fichier>}. Exemples d'utilisation :

\begin{lstlisting}
# Python
import pandas as pd
df = pd.read_csv("/work/data/population_sample.csv")
print(df.describe())
\end{lstlisting}

\begin{lstlisting}
# R
df <- read.csv("/work/data/population_sample.csv")
cat("Lignes :", nrow(df), "\n")
print(summary(df))
\end{lstlisting}

\subsection{Collecte et validation des fichiers de sortie}

\begin{table}[H]
\centering
\caption{Workflow de validation des outputs}
\begin{tabularx}{\textwidth}{@{}clX@{}}
\toprule
Etape & Acteur & Action \\
\midrule
1 & Script         & Ecrit ses fichiers dans \texttt{/work/outputs/} \\
2 & Worker         & Scanne le repertoire local apres execution \\
3 & Worker         & Insere dans \texttt{submission\_outputs} (statut PENDING\_VALIDATION) \\
4 & Administrateur & Approuve ou rejette chaque fichier \\
5 & Utilisateur    & Telecharge si APPROVED \\
\bottomrule
\end{tabularx}
\end{table}

\begin{lstlisting}
-- Table des fichiers de sortie
CREATE TABLE submission_outputs (
    id            SERIAL PRIMARY KEY,
    submission_id INTEGER REFERENCES submissions(id) ON DELETE CASCADE,
    filename      VARCHAR(255) NOT NULL,
    file_path     TEXT NOT NULL,
    created_at    TIMESTAMPTZ DEFAULT NOW()
);
\end{lstlisting}

\subsubsection*{Correction du bug chemin local / chemin hote}

\textbf{Probleme :} le worker (dans son conteneur) ne peut pas utiliser le chemin
Windows de l'hote pour \texttt{os.listdir}. \textbf{Solution :} deux chemins distincts.

\begin{lstlisting}
# Chemin pour Docker Desktop (-v mount)
host_outputs_dir = os.path.join(HOST_STORAGE_PATH, "outputs", str(sid))
# ex: /host_mnt/c/Users/.../storage/outputs/42

# Chemin accessible depuis le worker (volume docker-compose)
local_outputs_dir = os.path.join(STORAGE_PATH, "outputs", str(sid))
# ex: /storage/outputs/42

execute_script(file_path, microdata_path,
    outputs_dir=host_outputs_dir,        # pour -v
    local_outputs_dir=local_outputs_dir) # pour makedirs + listdir

collect_outputs(submission_id, local_outputs_dir, conn)
\end{lstlisting}

% ============================================================
\section{Issue \#7 -- Navigateur securise Brave}

\noindent\colorbox{myblue!10}{\parbox{\linewidth}{\vspace{4pt}
\textbf{Objectif :} restreindre l'acces a la plateforme au seul navigateur Brave
avec l'extension Brave Security Shield (MV3) installee, et appliquer les
protections cote client.
\vspace{4pt}}}

\vspace{6pt}

\subsection{Architecture de l'extension (Manifest V3)}

\begin{table}[H]
\centering
\caption{Fichiers de l'extension}
\begin{tabularx}{\textwidth}{@{}>{\ttfamily}lX@{}}
\toprule
\normalfont Fichier & Role \\
\midrule
manifest.json      & Declaration MV3, permissions, hotes autorises \\
content\_script.js & Protections injectees dans chaque page \\
background.js      & Service worker -- violations, gestion identite \\
popup.html/js      & Interface admin : toggles, watermark, journal \\
screen\_block.css  & Overlay CSS de blocage visuel \\
auth\_bridge.js    & API JS pour transmettre l'identite Flask vers BSS \\
\bottomrule
\end{tabularx}
\end{table}

\begin{table}[H]
\centering
\caption{Protections implementees}
\begin{tabularx}{\textwidth}{@{}lXc@{}}
\toprule
Protection & Mecanisme & Actif \\
\midrule
Copier/Coller    & Interception copy, cut, paste + Clipboard API      & \ok \\
Capture ecran    & Detection PrintScreen + overlay CSS + @media print & \ok \\
Partage ecran    & Remplacement de \texttt{getDisplayMedia()}          & \ok \\
DevTools         & 7 techniques : taille fenetre, F12, debugger...     & \ok \\
Filigrane        & SVG diagonal avec UID + date, configurable          & \ok \\
Journal          & chrome.storage.local, 200 entrees max               & \ok \\
\bottomrule
\end{tabularx}
\end{table}

\subsection{Integration Flask}

\subsubsection*{Detection par cookie de presence}

L'extension pose le cookie \texttt{bss\_active=1} (duree 24h) a chaque visite.
Flask verifie ce cookie sur \textbf{toutes les routes} via \texttt{browser\_required} :

\begin{lstlisting}
# backend/app/utils/authz.py
def browser_required(fn):
    @wraps(fn)
    def wrapper(*args, **kwargs):
        if not request.cookies.get("bss_active"):
            # HTTP 403 + page d'installation
            return render_template("browser_required.html"), 403
        return fn(*args, **kwargs)
    return wrapper
\end{lstlisting}

Si le cookie est absent : HTTP 403 + page d'installation de l'extension.
Si present : acces autorise normalement.

\subsubsection*{Filigrane -- liaison Flask / Extension}

Flask injecte l'identite de l'utilisateur dans la page via Jinja2 :

\begin{lstlisting}
<!-- En bas de chaque template protege -->
<script>
  var BSS_USER_ID  = {{ session.get('user_id') }};
  var BSS_USERNAME = {{ session.get('username') | tojson }};
</script>
<script type="module" src="/static/js/bss_init.js"></script>
\end{lstlisting}

\begin{lstlisting}
// static/js/bss_init.js
import { initBSSAuth } from "./auth_bridge.js";
if (typeof BSS_USER_ID !== "undefined" && BSS_USER_ID) {
    initBSSAuth({ uid: String(BSS_USER_ID), displayName: BSS_USERNAME || null });
}
\end{lstlisting}

\subsection{Installation de l'extension}

\begin{enumerate}
  \item Telecharger \textbf{Brave Browser} sur \url{https://brave.com}
  \item Ouvrir \texttt{brave://extensions} et activer le mode developpeur
  \item Cliquer sur <<\,Charger l'extension non empaquetee\,>>
  \item Selectionner le dossier \texttt{breavescripts/} du projet
  \item Copier l'Extension ID et le saisir dans \texttt{auth\_bridge.js} (ligne 37)
  \item Acceder a la plateforme : le filigrane et les protections sont actifs
\end{enumerate}

L'extension est installee et configuree avec l'ID reel :

\begin{lstlisting}
// breavescripts/auth_bridge.js  et  backend/static/js/auth_bridge.js
const BSS_EXTENSION_ID = "dikpdjlignemlnikaegcblghblbejfjd";
\end{lstlisting}

\subsection{Tests et validation en conditions reelles}

Les tests ont ete effectues dans Brave avec l'extension chargee et la plateforme
en cours d'execution (\texttt{http://localhost:5000}).

\begin{table}[H]
\centering
\caption{Resultats des tests de l'Issue \#7}
\begin{tabularx}{\textwidth}{@{}lXc@{}}
\toprule
\textbf{Test} & \textbf{Comportement observe} & \textbf{Statut} \\
\midrule
Acces depuis Chrome/Firefox
  & Page HTTP 403 affichee, acces refuse
  & \ok \\
Acces depuis Brave sans extension
  & Page HTTP 403 affichee, acces refuse
  & \ok \\
Acces depuis Brave avec extension
  & Plateforme accessible normalement
  & \ok \\
Filigrane sans connexion
  & Texte <<\,CONFIDENTIEL\,>> diagonal visible sur toutes les pages
  & \ok \\
Filigrane apres connexion (user: abdel)
  & Filigrane affiche <<\,abdel\,>> avec date en diagonal
  & \ok \\
Touche PrintScreen
  & Overlay noir 1,8 secondes, violation enregistree
  & \ok \\
Touche F12 / DevTools
  & Overlay <<\,ACCES REFUSE\,>> avec UID et timestamp
  & \ok \\
Ctrl+C / Ctrl+V
  & Toast rouge <<\,Copier/Coller desactive\,>>
  & \ok \\
Clic droit
  & Menu contextuel desactive
  & \ok \\
\bottomrule
\end{tabularx}
\end{table}

\noindent\colorbox{myblue!10}{\parbox{\linewidth}{\vspace{4pt}
\textbf{Limite documentee :} les outils de capture niveau OS (Win+Shift+S, OBS, Snagit)
ne peuvent pas etre bloques par une extension navigateur -- cela necessite un agent
systeme (C++). Cette limite est explicitement mentionnee dans la documentation
du prof (\texttt{breavescripts/readme.md}).
\vspace{4pt}}}

% ============================================================
\section{Benchmark du sandbox}

\subsection{Protocole}

\begin{table}[H]
\centering
\caption{Conditions du benchmark}
\begin{tabularx}{\textwidth}{@{}lX@{}}
\toprule
Parametre & Valeur \\
\midrule
Image Docker   & \texttt{portwatch-python-runner:1.0} \\
Limites        & 512 Mo RAM, 1 CPU, reseau desactive (\texttt{--network none}) \\
Timeout        & 60 secondes par soumission \\
Methode        & 5 executions par scenario, mesure du temps total docker run \\
Environnement  & Windows 11 Pro, Docker Desktop 4.x (WSL2) \\
Total          & 8 scenarios x 5 runs = \textbf{40 executions} \\
\bottomrule
\end{tabularx}
\end{table}

\subsection{Resultats}

\begin{table}[H]
\centering
\caption{Temps d'execution par scenario (ms)}
\begin{tabularx}{\textwidth}{@{}Xrrrrrr@{}}
\toprule
Scenario & R1 & R2 & R3 & R4 & R5 & Moy. \\
\midrule
\multicolumn{7}{@{}l}{\textit{Python}} \\
\quad Hello world             & 1611 & 2070 & 1417 & 1689 & 1379 & \textbf{1633} \\
\quad Pandas describe (1000L) & 1347 & 1502 & 1351 & 1271 & 1361 & \textbf{1366} \\
\quad Sklearn LinearRegression& 1395 & 1514 & 1253 & 1484 & 1427 & \textbf{1414} \\
\quad Lecture microdata CSV   & 1293 & 1593 & 1395 & 1545 & 1737 & \textbf{1512} \\
\quad Ecriture output CSV     & 1570 & 1601 & 1501 & 1344 & 1398 & \textbf{1482} \\
\midrule
\multicolumn{7}{@{}l}{\textit{R}} \\
\quad Hello world             & 1613 & 1613 & 1375 & 1423 & 1344 & \textbf{1473} \\
\quad rnorm(1000) + stats     & 1743 & 1483 & 1339 & 1600 & 1484 & \textbf{1529} \\
\quad Lecture microdata CSV   & 1424 & 1759 & 1631 & 1464 & 1332 & \textbf{1522} \\
\bottomrule
\end{tabularx}
\end{table}

\subsection{Analyse}

\textbf{Demarrage dominant :} la latence (1,3--1,7 s) est causee par le demarrage
du conteneur Docker. Le calcul lui-meme (pandas, sklearn, R stats) prend quelques
millisecondes -- le sandbox isole n'introduit pas de penalite de calcul significative.

\textbf{Python vs R -- performances equivalentes :} la difference moyenne est de
\textbf{27 ms} (Python 1481 ms, R 1508 ms), en dessous de la variabilite naturelle
(ecart-type ~150 ms). Les deux langages sont interchangeables.

\textbf{Overhead microdata/outputs negligeable :} ajouter un CSV en entree ou en
sortie represente moins de 10\% de la latence totale.

\begin{table}[H]
\centering
\caption{Verification des contraintes}
\begin{tabularx}{\textwidth}{@{}lXc@{}}
\toprule
Contrainte & Resultat & Statut \\
\midrule
Isolation reseau  & Aucun paquet sortant detecte        & \ok \\
Memoire 512 Mo    & Jamais atteinte                     & \ok \\
CPU 1 core        & Respecte                            & \ok \\
Timeout 60 s      & Moyenne 1,5 s (marge x40)           & \ok \\
Python            & pandas, numpy, sklearn disponibles  & \ok \\
R                 & stats, utils, base disponibles      & \ok \\
Microdata         & /work/data/ fonctionnel             & \ok \\
Outputs           & /work/outputs/ fonctionnel          & \ok \\
Taux erreur       & 0 erreur / 40 executions            & \ok \\
\bottomrule
\end{tabularx}
\end{table}

% ============================================================
\section{Conclusion}

\noindent\colorbox{green!8}{\parbox{\linewidth}{\vspace{6pt}
\textbf{Les deux issues ont ete integralement implementees et mergees sur \texttt{main} :}
\begin{itemize}
  \item \textbf{Issue \#6} -- Support R + selecteur microdata + collecte/validation des
        outputs avec workflow administrateur complet
  \item \textbf{Issue \#7} -- Extension Brave Security Shield : enforcement HTTP 403,
        filigrane utilisateur, 6 protections anti-exfiltration actives
\end{itemize}
\vspace{6pt}}}

\vspace{8pt}
Le benchmark confirme un cout fixe de \textbf{1,5 s} de demarrage sandbox, largement
acceptable pour une plateforme scientifique. Python et R sont a performances equivalentes,
et aucune des 40 executions n'a produit d'erreur.

\vspace{1.5cm}
\begin{flushright}
  \textbf{Abdelmoumen Elimrany}\\
  \textcolor{mygray}{22 avril 2026}\\
  \small\url{github.com/ElimranyAbdelmoumen/PET_Projet}
\end{flushright}

% ============================================================
\newpage
\section*{Annexe -- Structure du projet}
\addcontentsline{toc}{section}{Annexe -- Structure du projet}

\begin{lstlisting}
PET_Projet/
|-- backend/
|   |-- app/
|   |   |-- __init__.py        # Factory Flask, blueprint web
|   |   |-- routes/web.py      # Toutes les routes HTML
|   |   |-- static/js/         # auth_bridge.js, bss_init.js
|   |   |-- templates/         # login, submit, admin_*, user_view_*
|   |   +-- utils/authz.py     # login/admin/browser_required
|   +-- requirements.txt
|-- benchmark/
|   |-- run_benchmarks.sh      # 40 runs
|   +-- results.txt
|-- breavescripts/             # Extension Brave Security Shield MV3
|   |-- manifest.json
|   |-- content_script.js
|   |-- background.js
|   +-- auth_bridge.js
|-- db/init.sql                # Schema PostgreSQL
|-- docker/docker-compose.yml  # Backend + DB + Worker
|-- rapport/
|   |-- rapport.tex
|   +-- rapport.pdf
|-- runner/
|   |-- Dockerfile             # Sandbox Python + R
|   +-- run_script.py
|-- runner_worker/worker.py    # Boucle asynchrone
+-- storage/microdata/
\end{lstlisting}

\end{document}
